\documentclass[10pt]{article}
\usepackage[a4paper,margin=1in]{geometry}
\usepackage[T1]{fontenc}
\usepackage{lmodern}
\usepackage{amsmath,amssymb,mathtools,bm}
\usepackage{algorithm}
\usepackage[noend]{algpseudocode}
\usepackage{booktabs}
\usepackage{array}
\usepackage{microtype}
\usepackage{hyperref}

\algrenewcommand\algorithmicrequire{\textbf{Input:}}
\algrenewcommand\algorithmicensure{\textbf{Output:}}

\newcommand{\E}{\mathbb{E}}
\newcommand{\clip}{\operatorname{clip}}
\newcommand{\uid}{\operatorname{uid}}
\newcommand{\trajid}{\operatorname{traj\_uid}}
\newcommand{\stepidx}{\operatorname{step\_idx}}

\title{StepPPO-v2: GiGPO-Aligned Step-wise PPO with a Shared Value Head}
\author{}
\date{}

\begin{document}
\maketitle
\vspace{-1.5em}

\begin{abstract}
StepPPO-v2 is a shared-backbone actor-critic method for multi-step agentic reinforcement learning. An episode contains multiple environment interaction steps, and each step produces one complete language-model action. StepPPO-v2 defines credit assignment over environment steps rather than response tokens. The method keeps the GiGPO-style group-relative episode signal, adds a learned step-wise temporal correction from a shared value head, and optimizes the language model with the standard token-level PPO ratio. The value head is attached to the actor backbone; no separate critic model is maintained.
\end{abstract}

\section{Agentic Rollout Formulation}
We consider tasks where each sampled episode contains a sequence of agent-environment interaction steps. At environment step $t$, the agent receives context $s_t$, generates an action $a_t$, receives scalar reward $r_t$, and transitions to $s_{t+1}$. A trajectory is
\begin{equation}
\tau=(s_0,a_0,r_0,s_1,a_1,r_1,\ldots,s_{T-1},a_{T-1},r_{T-1}).
\end{equation}
The action $a_t$ is a complete response generated by the language model,
\begin{equation}
a_t=(y_{t,1},y_{t,2},\ldots,y_{t,L_t}),
\end{equation}
but StepPPO-v2 treats the whole response as a single environment-level action. Thus value estimation, temporal-difference errors, returns, and GAE are computed on the environment-step axis.

Each collected step sample $i$ stores three identifiers:
\begin{itemize}
\item $u_i$: the group id, corresponding to \texttt{uid}. It identifies repeated rollouts from the same initial task or prompt.
\item $\tau_i$: the trajectory id, corresponding to \texttt{traj\_uid}. It identifies one sampled episode.
\item $t_i$: the within-trajectory step index, corresponding to \texttt{step\_idx}. It orders samples inside the same trajectory.
\end{itemize}
The temporal transition for step-wise GAE is defined by grouping samples with the same $\tau_i$ and sorting them by $t_i$.

\section{Shared-Backbone Value Head}
Let $\pi_\theta$ be the actor language model. StepPPO-v2 attaches a small value head $v_\psi$ to the actor backbone. Given the tokenized context-action sequence for a step, the actor produces final-layer hidden states
\begin{equation}
H_i=f_\theta(s_i,a_i)\in\mathbb{R}^{n_i\times d},
\end{equation}
where $n_i$ is the full sequence length and $d$ is the hidden size.

StepPPO-v2 uses the pre-action boundary hidden state to produce a strict state-value estimate. If the generated action has response length $L_i$, the value index is
\begin{equation}
b_i=n_i-L_i-1.
\end{equation}
The predicted state value is
\begin{equation}
V_{\theta,\psi}(s_i)=v_\psi(H_i[b_i]).
\end{equation}
Because the backbone is causal, $H_i[b_i]$ cannot attend to generated action tokens. The value is therefore a baseline for the state before action generation.

The value head is an MLP regression head:
\begin{equation}
v_\psi(h)=W_2\phi(W_1h+b_1)+b_2,
\end{equation}
where $\phi$ is GELU. StepPPO-v2 does not detach the value head from the actor backbone:
\begin{equation}
\nabla_\theta \mathcal L_v \neq 0.
\end{equation}
Thus the value loss updates both the value-head parameters $\psi$ and the shared actor backbone $\theta$. To reduce interference from value regression, StepPPO-v2 uses a smaller value-loss coefficient.

\section{GiGPO-Aligned Episode Advantage}
For each step sample $i$, let $S_i$ denote the episode-level score. In implementation, $S_i$ is read from the collected episode return field, so every valid step in the same sampled episode receives the same episode score. Equivalently,
\begin{equation}
S_i=\sum_{t=0}^{T(\tau_i)-1}r_{\tau_i,t}.
\end{equation}

StepPPO-v2 aligns the episode advantage with the WebShop GiGPO baseline, which uses \texttt{mean\_norm}. For each group id $u$, define the row-level group set
\begin{equation}
\mathcal I_u=\{i\in\{1,\ldots,N\}:u_i=u\},
\qquad
n_u=|\mathcal I_u|.
\end{equation}
The group mean is
\begin{equation}
\mu_u=\frac{1}{n_u}\sum_{j\in\mathcal I_u}S_j.
\end{equation}
The episode advantage is the group-centered score
\begin{equation}
A_i^{\mathrm{epi}}=S_i-\mu_{u_i}.
\end{equation}
No additional batch-level whitening is applied to $A_i^{\mathrm{epi}}$. This keeps the episode term on the same normalization semantics as GiGPO.

\section{Step-wise Critic Return and Advantage}
Before PPO updates, StepPPO-v2 evaluates the behavior policy once to obtain both old token log probabilities and old state values:
\begin{equation}
\log\pi_{\theta_{\mathrm{old}}}(a_i\mid s_i),
\qquad
V_i^{\mathrm{old}}=V_{\theta_{\mathrm{old}},\psi_{\mathrm{old}}}(s_i).
\end{equation}
For a step sample $i=(\tau,t)$, let $i^+=(\tau,t+1)$ be the next step from the same trajectory if it exists. Let $d_i=1$ if $i$ is the final valid step in the trajectory and $d_i=0$ otherwise. The step-wise TD residual is
\begin{equation}
\delta_i=r_i+\gamma(1-d_i)V_{i^+}^{\mathrm{old}}-V_i^{\mathrm{old}}.
\end{equation}
The raw step advantage is computed by GAE on the environment-step axis:
\begin{equation}
A_i^{\mathrm{step,raw}}
=
\delta_i+
\gamma\lambda(1-d_i)A_{i^+}^{\mathrm{step,raw}}.
\end{equation}
The value target is
\begin{equation}
R_i^{\mathrm{step}}=A_i^{\mathrm{step,raw}}+V_i^{\mathrm{old}}.
\end{equation}

StepPPO-v2 standardizes the step advantage over all valid step samples in the current training batch:
\begin{equation}
\widehat A_i^{\mathrm{step}}
=
\frac{A_i^{\mathrm{step,raw}}-\mu_{\mathcal B}^{\mathrm{step}}}
{\sigma_{\mathcal B}^{\mathrm{step}}+\epsilon},
\end{equation}
where
\begin{equation}
\mu_{\mathcal B}^{\mathrm{step}}
=\frac{1}{N}\sum_{i=1}^{N}A_i^{\mathrm{step,raw}},
\qquad
\sigma_{\mathcal B}^{\mathrm{step}}
=\sqrt{\frac{1}{N}\sum_{i=1}^{N}(A_i^{\mathrm{step,raw}}-\mu_{\mathcal B}^{\mathrm{step}})^2}.
\end{equation}
This standardization controls the scale of the learned critic correction before it is combined with the GiGPO-aligned episode term.

\section{Final Advantage}
The final StepPPO-v2 policy advantage is
\begin{equation}
A_i^{\mathrm{v2}}=A_i^{\mathrm{epi}}+w\widehat A_i^{\mathrm{step}},
\end{equation}
with fixed coefficient
\begin{equation}
w=0.5.
\end{equation}
The combined advantage is not divided by $1+w$ and is not whitened after combination. It is broadcast to valid response tokens:
\begin{equation}
A_{i,\ell}^{\mathrm{v2}}=A_i^{\mathrm{v2}}m_{i,\ell},
\end{equation}
where $m_{i,\ell}$ is the response or loss mask. The step return is also broadcast for compatibility with the existing token-shaped value-loss path:
\begin{equation}
R_{i,\ell}^{\mathrm{step}}=R_i^{\mathrm{step}}m_{i,\ell}.
\end{equation}

\section{Optimization Objective}
For each response token, define the PPO ratio
\begin{equation}
\rho_{i,\ell}(\theta)=
\exp\left(
\log\pi_\theta(y_{i,\ell}\mid s_i,y_{i,<\ell})
-
\log\pi_{\theta_{\mathrm{old}}}(y_{i,\ell}\mid s_i,y_{i,<\ell})
\right).
\end{equation}
The clipped policy loss is
\begin{equation}
\mathcal L_{\mathrm{pg}}(\theta)
=-\frac{1}{\sum_{i,\ell}m_{i,\ell}}
\sum_{i=1}^{N}\sum_{\ell=1}^{L}
m_{i,\ell}
\min\left(
\rho_{i,\ell}(\theta)A_i^{\mathrm{v2}},
\clip(\rho_{i,\ell}(\theta),1-\epsilon_p,1+\epsilon_p)A_i^{\mathrm{v2}}
\right).
\end{equation}

For value learning, let
\begin{equation}
V_i=V_{\theta,\psi}(s_i),
\qquad
\widetilde V_i=V_i^{\mathrm{old}}+
\clip(V_i-V_i^{\mathrm{old}},-\epsilon_v,\epsilon_v).
\end{equation}
The clipped value loss is
\begin{equation}
\mathcal L_v(\theta,\psi)
=
\frac{1}{2N}\sum_{i=1}^{N}
\max\left(
(V_i-R_i^{\mathrm{step}})^2,
(\widetilde V_i-R_i^{\mathrm{step}})^2
\right).
\end{equation}
The total update loss is
\begin{equation}
\mathcal L(\theta,
\psi)=
\mathcal L_{\mathrm{pg}}(\theta)
+c_v\mathcal L_v(\theta,\psi)
+c_{\mathrm{kl}}\mathcal L_{\mathrm{kl}}(\theta)
-c_H\mathcal H(\theta).
\end{equation}
StepPPO-v2 uses
\begin{equation}
c_v=0.03,
\qquad
\epsilon_v=0.5.
\end{equation}

\section{Algorithm}
\begin{algorithm}[h]
\caption{StepPPO-v2 training iteration}
\begin{algorithmic}[1]
\Require Actor $\pi_\theta$ with value head $v_\psi$; grouped prompts; rollout repeat count $K$; $\gamma$, $\lambda$, $w=0.5$, $c_v=0.03$.
\Ensure Updated actor and value head.
\State Sample initial tasks and repeat each task $K$ times.
\State Roll out multi-step trajectories with the behavior policy $\pi_{\theta_{\mathrm{old}}}$.
\State Store each valid action step as $(u_i,\tau_i,t_i,s_i,a_i,r_i,S_i,m_i)$.
\State Run one old-policy forward pass to compute old token log probabilities and old state values $V_i^{\mathrm{old}}$.
\State Compute $A_i^{\mathrm{epi}}=S_i-\mu_{u_i}$ within each \texttt{uid} group.
\State Group by \texttt{traj\_uid}, sort by \texttt{step\_idx}, and compute $A_i^{\mathrm{step,raw}}$ and $R_i^{\mathrm{step}}$.
\State Standardize $A_i^{\mathrm{step,raw}}$ over the batch to obtain $\widehat A_i^{\mathrm{step}}$.
\State Form $A_i^{\mathrm{v2}}=A_i^{\mathrm{epi}}+0.5\widehat A_i^{\mathrm{step}}$.
\State Broadcast $A_i^{\mathrm{v2}}$ and $R_i^{\mathrm{step}}$ to response-token shape with the loss mask.
\For{each PPO epoch and minibatch}
  \State Recompute current token log probabilities and current values $V_i$.
  \State Minimize $\mathcal L_{\mathrm{pg}}+0.03\mathcal L_v+c_{\mathrm{kl}}\mathcal L_{\mathrm{kl}}-c_H\mathcal H$.
\EndFor
\end{algorithmic}
\end{algorithm}

\section{Implementation Configuration}
The StepPPO-v2 method is implemented with the shared value-head actor path. The deterministic paper configuration is:
\begin{verbatim}
algorithm.step_ppo.step_advantage_w=0.5
algorithm.step_ppo.episode_mode=mean_norm
algorithm.step_ppo.normalize_episode_advantage=False
algorithm.step_ppo.normalize_step_advantage=True
algorithm.step_ppo.normalize_final_advantage=False
actor_rollout_ref.actor.value_head.detach_value_backbone=False
actor_rollout_ref.actor.value_head.value_loss_coef=0.03
actor_rollout_ref.actor.value_head.cliprange_value=0.5
actor_rollout_ref.actor.ppo_mini_batch_size=64
actor_rollout_ref.actor.ppo_micro_batch_size_per_gpu=8
actor_rollout_ref.rollout.log_prob_micro_batch_size_per_gpu=16
actor_rollout_ref.ref.log_prob_micro_batch_size_per_gpu=16
\end{verbatim}
The actor update micro-batch size is set to $8$ per GPU to match the GiGPO WebShop baseline. The rollout and reference log-probability micro-batches remain at $16$ per GPU, also matching the GiGPO script.

\section{Implementation Mapping}
\begin{table}[h]
\centering
\small
\begin{tabular}{>{\raggedright\arraybackslash}p{0.18\linewidth}>{\raggedright\arraybackslash}p{0.32\linewidth}>{\raggedright\arraybackslash}p{0.39\linewidth}}
\toprule
Component & Code path & StepPPO-v2 behavior \\
\midrule
Value head & \texttt{value\_head.py} & Attach MLP value head to actor hidden states \\
Actor path & \texttt{step\_ppo\_actor.py} & Compute token log probabilities and state values in actor forward \\
Old values & \texttt{compute\_log\_prob} path & Return old log probabilities and old state values together \\
Step GAE & \texttt{step\_ppo\_algos.py} & Group by \texttt{traj\_uid}, sort by \texttt{step\_idx} \\
Episode term & advantage function & Use \texttt{uid}-group GiGPO \texttt{mean\_norm} \\
Final advantage & advantage function & Use $A^{\mathrm{epi}}+0.5\widehat A^{\mathrm{step}}$ \\
Value loss & actor update & Use clipped value loss with $c_v=0.03$ and no detach \\
Micro-batch & training script & Set actor PPO micro-batch to $8$ per GPU \\
\bottomrule
\end{tabular}
\end{table}

\section{Differences from StepPPO-v1}
StepPPO-v2 differs from StepPPO-v1 only in the following deterministic design choices:
\begin{itemize}
\item \textbf{Episode advantage}: v1 uses GRPO-style episode advantage followed by extra batch-level whitening. v2 uses GiGPO-aligned \texttt{uid}-group \texttt{mean\_norm} and does not apply additional batch-level whitening to the episode term.
\item \textbf{Final advantage}: v1 uses $(\widehat A^{\mathrm{epi}}+w\widehat A^{\mathrm{step}})/(1+w)$. v2 uses direct addition, $A^{\mathrm{epi}}+w\widehat A^{\mathrm{step}}$.
\item \textbf{Step weight}: v1 uses $w=1.0$. v2 uses $w=0.5$.
\item \textbf{Value loss coefficient}: v1 uses $c_v=0.1$. v2 uses $c_v=0.03$.
\item \textbf{Value backbone gradient}: v1 already allows value-loss gradients to update the shared backbone. v2 keeps this choice with \texttt{detach\_value\_backbone=False}.
\item \textbf{Micro-batch}: v1 WebShop scripts used actor PPO micro-batch size $4$ per GPU. v2 aligns with GiGPO and uses actor PPO micro-batch size $8$ per GPU.
\item \textbf{Old-value collection}: v2 keeps the single-forward path that returns old token log probabilities and old state values together when the value head is enabled.
\end{itemize}

\end{document}
